%% ──────────────────────────────────────────────────────────────────
%%  Diabetic Retinopathy — Lab 5 Report (v3)
%%  Computer Vision, April 2026
%% ──────────────────────────────────────────────────────────────────
\documentclass[a4paper, 10pt, twocolumn]{article}

\usepackage[top=1.6cm, bottom=1.6cm, left=1.4cm, right=1.4cm, columnsep=0.65cm]{geometry}
\usepackage[utf8]{inputenc}
\usepackage[T1]{fontenc}
\usepackage{lmodern}
\usepackage{microtype}
\usepackage{booktabs}
\usepackage{multirow}
\usepackage[table]{xcolor}
\usepackage{titlesec}
\usepackage[hidelinks]{hyperref}
\usepackage{enumitem}
\usepackage{amsmath}
\usepackage[labelfont=bf, font=small, skip=3pt]{caption}
\usepackage{tabularx}
\usepackage{array}
\usepackage{graphicx}

%% ── Colour palette ────────────────────────────────────────────────
\definecolor{accent}{RGB}{0, 64, 128}
\definecolor{rowbest}{RGB}{220, 240, 220}
\definecolor{rowbad} {RGB}{255, 230, 230}
\definecolor{rownorm}{RGB}{250, 250, 252}
\definecolor{rulecolor}{RGB}{0, 64, 128}

%% ── Section headings ─────────────────────────────────────────────
\titleformat{\section}
  {\normalfont\bfseries\small\color{accent}}
  {}{0em}{}[{\color{rulecolor}\vspace{-1pt}\hrule height 0.6pt\vspace{1pt}}]
\titlespacing{\section}{0pt}{4pt}{1pt}

\titleformat{\subsection}
  {\normalfont\small\bfseries}
  {}{0em}{}
\titlespacing{\subsection}{0pt}{2pt}{0pt}

%% ── Paragraph / list style ───────────────────────────────────────
\setlength{\parskip}{1pt}
\setlength{\parindent}{0pt}
\setlist[itemize]{topsep=0pt, itemsep=0pt, leftmargin=10pt,
                  label=\raisebox{0.3ex}{\tiny$\bullet$}, parsep=0pt}

%% ── Captions ─────────────────────────────────────────────────────
\setlength{\abovecaptionskip}{3pt}
\setlength{\belowcaptionskip}{0pt}

%% ══════════════════════════════════════════════════════════════════
\begin{document}

%% ── Title spanning both columns ──────────────────────────────────
\twocolumn[{%
\begin{center}
  {\large\bfseries\color{accent}
    Automatic Diagnosis of Diabetic Retinopathy\\[2pt]
    via Custom CNN Architecture and Transfer Learning}\\[3pt]
  {\small
    Maria Mu\~noz P\'erez \enspace$\cdot$\enspace
    Cristina Sobrino Fern\'andez-Miranda\enspace$\cdot$\enspace
    Computer Vision --- Lab 5 \enspace$\cdot$\enspace April 2026}\\[2pt]
  {\color{rulecolor}\hrule height 0.8pt}
  \vspace{4pt}
\end{center}
}]

%% ── 1. INTRODUCTION & DATA ──────────────────────────────────────
\section{Problem and Data}

We address binary classification of Diabetic Retinopathy (DR/No-DR) from
fundus photographs (2\,000 train / 500 val / 1\,000 test, 73\%/27\%
class imbalance) under two competition tracks: a CNN trained from scratch
(\textbf{CUSTOM}) and a fine-tuned ImageNet backbone (\textbf{FINE-TUNING}).
Performance is measured by AUC on Codabench. Below we describe the
\emph{final} solution submitted in each track; experiments that did not
make it into the final pipeline are summarised in the ablation table on
page\,2.

\textbf{Shared preprocessing.} Right-eye images are horizontally mirrored
to unify orientation. \texttt{CropByEye} segments the retinal disc by
intensity threshold; Ben~Graham contrast normalisation~\cite{graham2015kaggle}
($4{\times}\text{img}-4{\times}\text{GaussianBlur}+128$, with circular
mask and a \texttt{uint8} dtype guard) enhances microaneurysms and vessel
edges. ImageNet mean/std normalisation closes the pipeline.

\textbf{Class imbalance.} Two complementary mechanisms are used together:
\texttt{BCEWithLogitsLoss} with \texttt{pos\_weight}\,$\approx$\,2.76
(gradient-level reweighting) and a \texttt{WeightedRandomSampler} with
inverse-frequency weights (batch-level rebalancing).

%% ── 2. CUSTOM MODEL ────────────────────────────────────────────
\section{CUSTOM: 3-Seed CustomNetV2 Ensemble}

\textbf{Architecture (CustomNetV2, $\sim$1\,M params).}
Five convolutional blocks of \texttt{Conv2d(3$\times$3, same-pad, no bias)
$\to$ BN $\to$ ReLU $\to$ MaxPool(2)} with channel progression
$3{\to}32{\to}64{\to}128{\to}256{\to}256$. Spatial path
$224{\to}112{\to}56{\to}28{\to}14{\to}7$. A \emph{Global Average Pooling}
layer replaces the original flatten (translation invariance, far fewer
parameters in the head), followed by
\texttt{Dropout(0.2) $\to$ Linear(256,128) $\to$ ReLU $\to$ Dropout(0.4)
$\to$ Linear(128,1)}.

\textbf{Training pipeline (224\,px).}
\texttt{CropByEye} $\to$ \texttt{BenGraham} $\to$ \texttt{Rescale(256)}
$\to$ random h/v flip $\to$ rotation\,$\pm$15$^\circ$ $\to$ ColorJitter
$\to$ \texttt{RandomCrop(224)} $\to$ \texttt{Normalize}. Validation/test
use \texttt{CenterCrop(224)} only.

\textbf{Optimisation.} \texttt{AdamW}~(lr=1e-3, wd=5e-3),
\texttt{CosineAnnealingLR}($T_{\max}{=}50$), 50 epochs, patience\,=\,10
on val AUC. Weighted sampler with \texttt{shuffle=False}.

\textbf{Final model: 3-seed ensemble.} Three independent runs with seeds
$\{42, 123, 456\}$ are uniformly averaged at inference. Diversity comes
from random initialisation, batch ordering and augmentation rather than
architectural complexity --- earlier attempts adding SE attention or
residual connections (CustomNetV3, SECustomNetV3) all overfit on the
2\,000-image set (Tab.\,\ref{tab:abl}). Each model is evaluated with
\textbf{4-pass TTA} (original, hflip, vflip, rot180) before averaging.
\textbf{Codabench: 0.757 pub / 0.780 priv.}

%% ── 3. FINE-TUNING MODEL ───────────────────────────────────────
\section{FINE-TUNING: 3-Backbone Weighted Ensemble at 380\,px}

\textbf{Backbones.} Three ImageNet-pretrained backbones with complementary
inductive biases are combined: \textbf{DenseNet121}~\cite{huang2017densely}
(dense feature reuse), \textbf{DenseNet169} (deeper dense reuse) and
\textbf{ResNeXt50-32x4d} (grouped convolutions). Each has a custom head
\texttt{Dropout(0.5) $\to$ Linear($d$,1)} where $d$ is the backbone's
feature dimension (1024 / 1664 / 2048).

\textbf{Two-stage progressive unfreezing.} Each backbone is trained
identically:
\begin{itemize}
  \item \emph{Stage 1.} Backbone frozen except for the deepest block + norm
    (DenseNet: \texttt{denseblock4+norm5}; ResNeXt: \texttt{layer4})
    plus the head. \texttt{AdamW} lr=3e-4, wd=1e-2,
    \texttt{CosineAnnealingLR}($T_{\max}{=}30$), 30 epochs, patience=7,
    label smoothing $\varepsilon{=}0.05$.
  \item \emph{Stage 2.} Stage-1 weights loaded; one block deeper additionally
    unfrozen (\texttt{denseblock3+transition3} or \texttt{layer3}).
    Per-group LRs: backbone 3e-5, head 1e-4. \texttt{LinearLR} 3-epoch
    warm-up (the warm-up prevents gradient spikes when freshly unfrozen,
    unadapted layers first see the fine-tuning LR) then
    \texttt{CosineAnnealingLR}($T_{\max}{=}12$), 15 epochs, patience=5.
    A gate skips Stage\,2 if Stage\,1 + TTA val AUC $<$ 0.752, and the
    Stage\,1 checkpoint is restored if Stage\,2 regresses.
\end{itemize}

\textbf{High-resolution pipeline (380\,px).} Microaneurysms (10--125\,$\mu$m)
span $\sim$2\,px at 224\,px, below the reliable detection limit of a
$3{\times}3$ filter; at 380\,px they reach 3--4\,px and are recoverable.
This is the primary driver of sensitivity for early DR
lesions~\cite{abramoff2018pivotal}. Pipeline: \texttt{CropByEye} $\to$
\texttt{BenGraham} $\to$ \texttt{Rescale(416)} $\to$ random hflip $\to$
\textbf{rotation\,$\pm$180$^\circ$} $\to$ ColorJitter $\to$
\texttt{RandomCrop(380)} $\to$ \texttt{Normalize}. The full $\pm$180$^\circ$
range is used because fundus cameras have no canonical orientation; the
previous $\pm$15$^\circ$ window discarded $\sim$92\% of valid rotations,
and full rotation is reported as the most impactful fundus
augmentation~\cite{graham2015kaggle}. Backbones are fully convolutional
(GAP head), so no architectural change is needed for the higher resolution;
batch sizes drop to 32/128 to fit the GPU.

\textbf{Inference \& ensemble.} Each backbone is scored with the same
4-pass TTA. A grid search over all 2/3/4-backbone combinations and
discrete weights $\{0.10,0.15,\dots,0.80\}$ on the 500-image val set
selects:
\[
\text{DenseNet121}{\times}0.35 + \text{DenseNet169}{\times}0.15
   + \text{ResNeXt50}{\times}0.50 .
\]
SE-ResNeXt50 (also trained) was \emph{discarded} by the search: weaker
individually and positively correlated with ResNeXt50, it dilutes the
blend. \textbf{Val AUC 0.867. Codabench: 0.835 pub / 0.826 priv.}

%% ── 4. SUMMARY ─────────────────────────────────────────────────
\section{Final Results}

\begin{tabular}{@{}lcc@{}}
\toprule
Track       & Public AUC & Private AUC \\
\midrule
CUSTOM (3-seed CustomNetV2) & 0.757 & \textbf{0.780} \\
FINE-TUNING (3-backbone, 380\,px)   & \textbf{0.835} & \textbf{0.826} \\
\bottomrule
\end{tabular}


%% ══════════════════════════════════════════════════════════════════
%% PAGE 2 — ABLATIONS, NEGATIVE RESULTS & REFERENCES
%% ══════════════════════════════════════════════════════════════════
\newpage
\onecolumn

\section*{\color{accent}Page 2 --- Ablations, Negative Results and References}
{\color{rulecolor}\hrule height 0.6pt}
\vspace{4pt}

%% ── ABLATIONS TABLE ──────────────────────────────────────────────
\captionof{table}{Changes that were tried during development but did
\emph{not} end up in the final solutions of either track. Each entry
records the change, the observed effect, and why it was rejected.
\textit{Cust} = CUSTOM track, \textit{FT} = FINE-TUNING track.}
\label{tab:abl}
\vspace{3pt}
\footnotesize
\setlength{\tabcolsep}{4pt}
\renewcommand{\arraystretch}{1.18}
\begin{tabularx}{\textwidth}{@{}p{0.5cm} p{4.0cm} p{3.4cm} X@{}}
\toprule
\textbf{Track} & \textbf{Change tried} & \textbf{Effect observed} & \textbf{Why discarded} \\
\midrule

\multicolumn{4}{@{}l}{\textit{\color{accent}Preprocessing}} \\
Both & CLAHE (LAB L-channel, clipLimit=2.0) before \texttt{CropByEye}
     & Cust $-$0.022 pub; FT priv $-$0.017
     & CLAHE alters the global intensity distribution; \texttt{CropByEye}'s
       fixed threshold (0.10) then mis-segments the eye disc. Permanently banned. \\

Both & BenGraham without \texttt{uint8} dtype guard
     & Catastrophic: Cust 0.541, FT 0.558 (both near random)
     & \texttt{np.clip(uint8,0,1)} collapsed every non-zero pixel to 1, so
       $4{\cdot}\text{img}-4{\cdot}\text{blur}+128 \approx 128$ everywhere.
       Models trained on uniform grey patches. Fixed by adding a dtype
       guard and a circular background mask. \\

Both & Pre-cached 224\,px tensors on disk (5$\times$ training speed-up)
     & Cust $-$0.030 pub; FT $-$0.028 pub
     & \texttt{CenterCrop} applied at cache-creation removes the per-epoch
       \texttt{RandomCrop} jitter, the dominant spatial regulariser on a
       2\,000-image set. Speed-up not worth the AUC loss. \\

\addlinespace[2pt]
\multicolumn{4}{@{}l}{\textit{\color{accent}Augmentation \& TTA}} \\
Both & 8-pass TTA (added rot90 / rot270 to the 4-pass set)
     & Both tracks regressed
     & BenGraham + rectangular crop is not invariant under 90$^\circ$
       rotation; rot90/rot270 images fall outside the training distribution
       and miscalibrate the mean. \\

FT   & 2-pass TTA (hflip only) instead of 4-pass
     & Lower AUC than 4-pass
     & vflip and rot180 add free inference-time diversity; the dataset has
       no strong vertical anatomical prior. \\

Cust & \texttt{TVRandomResizedCrop(scale 0.65--1.0)} replacing
       \texttt{RandomCrop(224)}, plus label smoothing $\varepsilon{=}0.05$
     & Marginal change vs.\ tuned baseline
     & Reverted: harder to reason about train/eval crop mismatch with
       fixed \texttt{CenterCrop}. Label smoothing kept only for FT,
       where pretrained logits benefit more~\cite{muller2019label}. \\

Cust & 380\,px input + $\pm$180$^\circ$ rotation on plain CustomNetV2
       (the same recipe that helped FT)
     & 3-seed val AUC = 0.530 (effectively random)
     & The custom model's $5{\times}\text{MaxPool}$ path is calibrated for
       a $7{\times}7$ feature map at 224\,px. At 380\,px GAP aggregates an
       $11{\times}11$ map, a different receptive-field regime; without
       ImageNet pretraining, training from scratch did not converge. \\

\addlinespace[2pt]
\multicolumn{4}{@{}l}{\textit{\color{accent}Custom architecture}} \\
Cust & SE blocks (reduction=8) on every CustomNetV2 block
     & SE-V2 alone reached only $\sim$0.52 val AUC; ensembles regressed
     & $\sim$22\,k extra params introduce channel attention that overfits
       on 2\,000 images without ImageNet-scale pretraining. \\

Cust & Residual connections (CustomNetV3:
       \texttt{Conv$\to$BN + 1$\times$1 skip $\to$ ReLU $\to$ Pool})
     & At best $+$0.002 over plain V2; ensembles below 3-seed V2
     & Skip connections matter when vanishing gradients are the bottleneck;
       a 5-block net trains fine without them. \\

Cust & Partial SE on blocks 3--5 only (reduction=16)
     & Val AUC $\sim$0.72
     & A late-stage-only attention does not resolve the underlying
       overfitting on 2\,000 images. \\

Cust & 4-architecture uniform-mean ensemble (V2 + SE-V2 + V3 + SE-V3, single seed)
     & 0.760 pub / 0.773 priv
     & Lower private than the plain-V2 3-seed ensemble (0.780). Architectural
       diversity introduced unstable members (SE-V2 $\sim$0.52 val).
       Stochastic-seed diversity proved more reliable. \\

\addlinespace[2pt]
\multicolumn{4}{@{}l}{\textit{\color{accent}Fine-tuning architecture \& training}} \\
FT   & SGD\,+\,StepLR (EfficientNet-B0 baseline)
     & Ceiling at 0.698 pub
     & Adaptive optimisers with decoupled weight decay (AdamW) + smooth
       cosine annealing transfer better on small medical-imaging datasets. \\

FT   & EfficientNet-B0 with three top blocks unfrozen
     & Best 0.770 pub / 0.744 priv (vs.\ DenseNet121 0.781/0.781 in
       development); train-val gap 0.178
     & Compound-scaled architecture overfits without large data;
       DenseNet's dense feature reuse acts as implicit regularisation. \\

FT   & DenseNet Stage\,3 unfreeze (\texttt{denseblock2 + transition2})
     & FT $-$0.013 pub vs.\ two-stage
     & Stages 1+2 already unfreeze $\sim$5.5\,M of 6.95\,M parameters;
       Stage\,3 over-parameterises without adding capacity. \\

FT   & EfficientNet-B4-NS (Noisy Student) as a 5th backbone at 380\,px
     & CUDA OOM (14.56\,GiB GPU saturated at batch=16, even with grad. checkpointing)
     & Hardware constraint, not a methodological one. Scientific motivation
       valid (APTOS-2019 winners used B3--B7) but incompatible with this setup. \\

FT   & SE-ResNeXt50 (timm) included in the final ensemble
     & Dropped by val grid search; every 3-backbone blend including it scored
       below the chosen DenseNet121 + DenseNet169 + ResNeXt50 blend
     & Single-model val AUC (0.834) below ResNeXt50 (0.855), and predictions
       are positively correlated with ResNeXt50 (same family) ---
       weaker \emph{and} less complementary. \\

\bottomrule
\end{tabularx}

\vspace{12pt}

%% ── KEY VALID SUBMISSIONS (CONTEXT) ──────────────────────────────
\captionof{table}{Key valid Codabench submissions illustrating the
trajectory of the development. Final models highlighted in green.}
\label{tab:traj}
\vspace{3pt}
\footnotesize
\setlength{\tabcolsep}{5pt}
\renewcommand{\arraystretch}{1.10}
\begin{tabularx}{\textwidth}{@{}lXcc@{}}
\toprule
\textbf{Notebook} & \textbf{Key change} &
  \textbf{Custom AUC} & \textbf{FT AUC} \\
 & & \scriptsize pub\,/\,priv & \scriptsize pub\,/\,priv \\
\midrule
\texttt{first\_attempt}        & EfficientNet-B0 (SGD) + CustomNet baseline
                                & 0.570\,/\,0.597 & 0.698\,/\,0.708 \\
\texttt{attempt\_with\_c}      & FT: AdamW + CosineAnnealingLR
                                & 0.561\,/\,0.579 & 0.751\,/\,0.734 \\
\texttt{second\_attempt\_tuned}& CustomNetV2 + WeightedRandomSampler + 2-stage FT
                                & 0.766\,/\,0.743 & 0.770\,/\,0.744 \\
\texttt{fifth\_attempt}        & BenGraham fixed + 4-pass TTA + label smoothing
                                & 0.770\,/\,0.770 & 0.757\,/\,0.755 \\
\texttt{attempt\_22}           & FT: ResNet50 (IMAGENET1K\_V2) single model
                                & 0.756\,/\,0.764 & 0.780\,/\,0.791 \\
\texttt{attempt\_21}           & FT: 4-backbone ensemble at 224\,px
                                & 0.757\,/\,0.780 & 0.806\,/\,0.807 \\
\rowcolor{rowbest}
\texttt{ninth\_attempt}        & \textbf{CUSTOM final} --- plain CustomNetV2, 3-seed ensemble
                                & \textbf{0.757\,/\,0.780} & --- \\
\rowcolor{rowbest}
\texttt{attempt\_24}           & \textbf{FT final} --- 380\,px + $\pm$180$^\circ$ rotation, 3-backbone weighted ensemble
                                & 0.760\,/\,0.773 & \textbf{0.835\,/\,0.826} \\
\bottomrule
\end{tabularx}

\vspace{12pt}

%% ── REFERENCES ───────────────────────────────────────────────────
\renewcommand{\refname}{\normalfont\bfseries\small\color{accent}References}
\begin{thebibliography}{9}
\small\setlength{\itemsep}{1pt}

\bibitem{graham2015kaggle}
Graham, B.\,(2015).
Kaggle diabetic retinopathy detection competition report.
\textit{University of Warwick}.

\bibitem{huang2017densely}
Huang, G., Liu, Z., van der Maaten, L., \& Weinberger, K.\,Q.\,(2017).
Densely connected convolutional networks.
\textit{Proc.\ CVPR 2017}, 4700--4708.

\bibitem{abramoff2018pivotal}
Abramoff, M.\,D., et al.\,(2018).
Pivotal trial of an autonomous AI-based diagnostic system for detection of
diabetic retinopathy in primary care offices.
\textit{npj Digital Medicine}, 1(1), 39.

\bibitem{loshchilov2019adamw}
Loshchilov, I., \& Hutter, F.\,(2019).
Decoupled weight decay regularization.
\textit{Proc.\ ICLR 2019}.

\bibitem{muller2019label}
M\"uller, R., Kornblith, S., \& Hinton, G.\,(2019).
When does label smoothing help?
\textit{Proc.\ NeurIPS 2019}, 4694--4703.

\bibitem{xie2020se}
Hu, J., Shen, L., \& Sun, G.\,(2018).
Squeeze-and-excitation networks.
\textit{Proc.\ CVPR 2018}, 7132--7141.

\end{thebibliography}

\end{document}
